\chapter{Overview}
\label{ch:overview}

\section{What the system does}

Big\_Agent turns a folder of engineering reports into something an engineer can
question. The pipeline is deliberately ordinary: parse, clean, chunk, embed,
retrieve, and answer only from what was retrieved. The deterministic half of
the system is the product; the language model is an optional layer on top of
it, and the application stays useful when that layer is switched off.

\section{Capabilities}

\begin{longtable}{L{0.26\textwidth} L{0.66\textwidth}}
\toprule
\textbf{Area} & \textbf{What is implemented} \\
\midrule
\endfirsthead
\toprule
\textbf{Area} & \textbf{What is implemented} \\
\midrule
\endhead
\bottomrule
\endfoot
Ingestion & Single-file and folder ingestion of PDF, DOCX and PPTX, duplicate
detection by file hash, selective OCR for scanned PDF pages. \\
Processing & Whitespace and Unicode normalization that preserves Turkish
characters, repeated header/footer removal, overlapping chunking. \\
Retrieval & Keyword search, semantic search over stored embeddings, hybrid
scoring, and an optional Haystack pipeline as a second retrieval
implementation. \\
Question answering & Source-grounded answers over one report, several reports,
or the report catalog; every answer carries the passages it was built from. \\
Catalog & Excel/catalog import, catalog-to-file matching over network shares,
preview and open workflows, and a graph overview of the corpus. \\
Review & Deterministic report-review rules with discipline profiles, human
confirm/dismiss decisions, per-rule precision, PDF review records, and revision
comparison. \\
Comparison & Pairwise and multi-report comparison with highlighted PDF
evidence. \\
Drafting & A report draft builder that writes a cover block and a body grounded
in retrieved passages, exportable as PDF. \\
CATIA skill & A mass/centre-of-gravity measurement skill driven from the chat
UI, backed by the \file{skill/catia-mass-cg.skill} package. \\
\end{longtable}

\section{Product variants}

One codebase serves three product identities, selected with \env{APP\_VARIANT}
(see \file{app/branding.py}).

\begin{longtable}{L{0.16\textwidth} L{0.18\textwidth} L{0.22\textwidth} L{0.32\textwidth}}
\toprule
\textbf{Variant} & \textbf{Display name} & \textbf{Data dir env} & \textbf{Notes} \\
\midrule
\endfirsthead
\bottomrule
\endfoot
\code{big\_agent} & SmartCAE AI & \env{BIG\_AGENT\_DATA\_DIR} & Default. The
only variant with a fully wired application workspace. \\
\code{raporhub} & RaporHub & \env{RAPORHUB\_DATA\_DIR} & Landing experience
served from \file{app/ui/raporhub\_landing}. \\
\code{repocto} & RepOcto & \env{REPOCTO\_DATA\_DIR} & Landing experience plus
the library browser (\route{POST /library/scan}). \\
\end{longtable}

The variant changes the API title, the browser icon, the palette and which
front end \route{/} serves. It does not change the retrieval or review engines.

\section{User interfaces}

\begin{description}
\item[\route{/}] The SmartCAE v2 workspace (\file{app/ui/smartcae\_v2}) for the
\code{big\_agent} variant, or the variant landing page otherwise.

\item[\route{/app}, \route{/legacy}] The earlier single-page interface in
\file{app/static}. Still served, still functional, and useful when debugging an
endpoint without the newer front end.

\item[\route{/docs}, \route{/redoc}] FastAPI's generated OpenAPI documentation
for the live instance. This manual describes intent and behaviour; \route{/docs}
is the authority on the exact schema of the version you are running.
\end{description}

\section{Design rules the project holds to}

These are worth stating up front, because most of the code makes more sense once
they are known.

\begin{enumerate}
\item \textbf{The deterministic retrieval system is the main system.} LLM
features are layered on top and must fail soft: an unavailable model degrades an
answer, it never breaks an endpoint.

\item \textbf{Defaults are offline.} \env{LLM\_ENABLED=false},
\env{LLM\_BACKEND=disabled}, \env{RERANKER\_ENABLED=false}. The application must
start with no model present at all.

\item \textbf{Embedding models and generation models are separate roles.} A
\code{Qwen3-Embedding-*} model produces retrieval vectors and nothing else; it
is never loaded with a causal-LM head.

\item \textbf{Turkish text is first-class.} Normalization folds Turkish
characters for matching but never destroys them in stored text.

\item \textbf{Parser, processing, database and service responsibilities stay
separate}, and interfaces stay modular so the embedding provider or vector
backend can be replaced.
\end{enumerate}

\section{Local data policy}

The following are intentionally not committed to the repository:

\begin{itemize}
\item \file{data/} --- the SQLite database, stored documents, job artefacts
\item \file{models/} --- local embedding model weights
\item report files: \file{.pdf}, \file{.docx}, \file{.pptx}, \file{.xlsx},
\file{.csv}
\item \file{.env} files
\end{itemize}

This keeps company documents, local databases and model weights out of version
control. Synthetic corpora produced by
\file{scripts/generate\_sample\_reports.py} are gitignored too; see
Chapter~\ref{ch:testing}.

\section{Version}

The running version is \code{APP\_VERSION} in \file{app/version.py}, returned by
\route{GET /health} and \route{GET /system/model-status} and used as the FastAPI
application version. This manual is written against \version.
